\chapter{Respaldos y restauración}
\label{cap:respaldos}

Este capítulo resume la operación en los tres modos. El procedimiento detallado, con la tabla de
errores de restauración, está en \cod{docs/respaldos.md} del repositorio.

\section{Cómo funcionan}

\begin{xltabular}{\linewidth}{@{}l L@{}}
	\toprule
	\textbf{Qué} & \textbf{Valor por defecto} \\
	\midrule
	\endhead
	Frecuencia & Cada 7 días (\texttt{BACKUP\_INTERVAL\_DAYS}) \\
	Hora & 06:00 UTC = 02:00 en Caracas (\texttt{BACKUP\_HOUR\_UTC}) \\
	Conservación & Se borran los de más de 70 días (\texttt{BACKUP\_KEEP\_DAYS}), pero siempre
	quedan los 10 más recientes (\texttt{BACKUP\_KEEP\_MIN}) \\
	Espacio mínimo & No empieza con menos de 512 MB libres o el doble del último respaldo
	(\texttt{BACKUP\_MIN\_FREE\_MB}) \\
	Respaldo perdido & Si el servidor estaba apagado, se hace en cuanto vuelve (si el último tiene
	más de 8 días). Al instalar, se hace uno de inmediato \\
	Bloqueo & Ninguno: instantánea consistente de MySQL; la aplicación sigue funcionando \\
	\bottomrule
\end{xltabular}

Cada respaldo es una carpeta \texttt{digaipe-<fecha UTC>} con \texttt{digaipe.sql.gz} (datos,
auditoría y procedimientos, sin triggers), \texttt{uploads/} (archivos cargados),
\texttt{MANIFEST} (fecha, versión de MySQL y cabeza de la cadena de auditoría) y
\texttt{SHA256SUMS}. Cada respaldo, y cada fallo, queda en la auditoría como
\texttt{backup.completed} o \texttt{backup.failed}.

\begin{aviso}
	No cambie el nombre de las carpetas de respaldo: la antigüedad se lee del nombre, y una carpeta
	con otro nombre nunca se cuenta ni se borra. Los respaldos contienen datos personales y hashes
	de contraseñas: guárdelos con el mismo cuidado que la base.
\end{aviso}

\section{Revisión semanal (5 minutos)}

\begin{enumerate}
	\item El último resultado debe ser \texttt{Backup finished.}:
	      \begin{itemize}
		      \item Docker: \texttt{docker compose logs --tail 20 backup-worker}
		      \item Linux: \texttt{journalctl -u digaipe-respaldo -n 20}
		      \item Windows: \texttt{Get-Content C:\bs digaipe\bs logs\bs respaldo.log -Tail 20}
	      \end{itemize}
	\item En Auditoría, filtrando por acción, debe haber un \texttt{backup.completed} de los
	      últimos 7 días y ningún \texttt{backup.failed} sin atender.
	\item Copie la carpeta del respaldo más reciente a un medio fuera del servidor.
\end{enumerate}

\section{Respaldo manual}

Por ejemplo, antes de una actualización:

\begin{xltabular}{\linewidth}{@{}l L@{}}
	\toprule
	\textbf{Modo} & \textbf{Comando} \\
	\midrule
	\endhead
	Docker & \texttt{docker compose --profile ops run --rm backup} \\
	Linux & \texttt{sudo systemd-run --wait --pipe -p User=digaipe-respaldo
	-p SupplementaryGroups=digaipe -p EnvironmentFile=/etc/digaipe/backup.env
	/usr/bin/bash /opt/digaipe/app/deploy/backup/backup.sh} \\
	Windows & \texttt{\& 'C:\bs Program Files\bs Git\bs bin\bs bash.exe' -c '. /c/digaipe/bin/entorno.sh
	\&\& bash /c/digaipe/app/deploy/backup/backup.sh'} \\
	\bottomrule
\end{xltabular}

\section{Antes de restaurar}

\textbf{Restaurar devuelve el sistema al momento del respaldo. Todo lo posterior se pierde} de
la base activa: admisiones, planillas, revisiones, usuarios y eventos de auditoría de ese
periodo. Restaure solo si la base está dañada o se perdió, o por decisión documentada ante un
incidente.

\begin{enumerate}
	\item Registre el incidente: qué pasó, quién autoriza, fecha y hora.
	\item Elija el respaldo más reciente anterior al problema y lea su \texttt{MANIFEST}.
	\item Avise a los usuarios: el sistema estará detenido unos minutos.
\end{enumerate}

\section{Restaurar}

En todos los modos, el orden es el mismo: detener la aplicación y los respaldos, restaurar,
\textbf{volver a ejecutar la instalación} (reconstruye los triggers de auditoría, que los
respaldos no traen), verificar la auditoría y arrancar. Reemplace
\texttt{digaipe-20261028T060000Z} por el respaldo elegido.

\so{Docker} (desde \texttt{deploy})
\begin{codigo}
docker compose stop app backup-worker
docker compose --profile ops run --rm -e RESTORE_CONFIRM=digaipe-20261028T060000Z restore /scripts/restore.sh digaipe-20261028T060000Z
docker compose run --rm setup
docker compose --profile ops run --rm audit verify
docker compose start app backup-worker
\end{codigo}

\so{Linux}
\begin{codigo}
sudo systemctl stop digaipe digaipe-respaldo
sudo /usr/local/sbin/digaipe-restaurar digaipe-20261028T060000Z digaipe-20261028T060000Z
sudo /usr/local/sbin/digaipe-setup
sudo /usr/local/sbin/digaipe-audit verify
sudo systemctl start digaipe digaipe-respaldo
\end{codigo}

\so{Windows}
\begin{codigo}
Stop-Service digaipe
Disable-ScheduledTask -TaskPath '\DIGAIPE\' -TaskName Respaldo
& 'C:\Program Files\Git\bin\bash.exe' C:/digaipe/bin/restaurar.sh digaipe-20261028T060000Z digaipe-20261028T060000Z
powershell -ExecutionPolicy Bypass -File C:\digaipe\bin\digaipe-setup.ps1
powershell -ExecutionPolicy Bypass -File C:\digaipe\bin\digaipe-audit.ps1 verify
Start-Service digaipe
Enable-ScheduledTask -TaskPath '\DIGAIPE\' -TaskName Respaldo
\end{codigo}

La restauración comprueba \texttt{SHA256SUMS}, hace antes un \textbf{respaldo de seguridad} del
estado actual (su nombre aparece en la salida: anótelo en el incidente y no lo borre), reemplaza
la base y los archivos, y comprueba la cabeza de la cadena de auditoría.

\begin{comprobar}
	Inicie sesión como administrador. En Auditoría deben aparecer \texttt{backup.restored} y,
	justo después, \texttt{schema.triggers\_repaired}. Fuera de una restauración, un
	\texttt{schema.triggers\_repaired} significa que alguien borró o cambió un trigger de
	auditoría: trátelo como incidente.
\end{comprobar}

\section{Recuperar en un servidor nuevo}

\begin{enumerate}
	\item Instale el servidor completo según su modo (capítulos~\ref{cap:docker}
	      a~\ref{cap:windows}) con los \textbf{mismos} archivos de configuración si los conserva,
	      pero \textbf{sin} ejecutar todavía la instalación ni arrancar la aplicación.
	      En Docker: \texttt{docker compose up -d db} (solo la base) y espere a que esté
	      \texttt{healthy}.
	\item Copie la carpeta del respaldo a la carpeta de respaldos del servidor nuevo.
	\item Siga la sección anterior desde el paso de restaurar, agregando
	      \texttt{SKIP\_SAFETY\_BACKUP=1} (no hay nada que respaldar): en Docker con
	      \texttt{-e SKIP\_SAFETY\_BACKUP=1}; en servidor físico, anteponiendo
	      \texttt{SKIP\_SAFETY\_BACKUP=1} al comando en Linux
	      (\texttt{sudo SKIP\_SAFETY\_BACKUP=1 /usr/local/sbin/digaipe-restaurar ...}) o agregando
	      \texttt{export SKIP\_SAFETY\_BACKUP=1} a \texttt{restaurar.sh} en Windows (y quitándolo
	      después).
	\item Instale la raíz HTTPS nueva en las PCs (capítulo~\ref{cap:certificado}), salvo que haya
	      conservado la carpeta de la autoridad.
\end{enumerate}

\section{Simulacro de restauración}
\label{sec:simulacro}

Un respaldo que nunca se restauró no está comprobado. Al menos una vez por periodo académico,
haga la sección anterior completa en un equipo aparte (nunca en el servidor de producción) con
el respaldo más reciente, y compruebe que puede iniciar sesión y ver las últimas admisiones.
Registre la fecha, el respaldo usado y el resultado.
